\documentclass[letter,11pt]{article}
\usepackage{hyperref}
\usepackage{graphicx}
\usepackage{fancyhdr}
\usepackage{multicol}
\pagestyle{fancy}
\usepackage[letterpaper, margin=1in]{geometry}
\geometry{letterpaper}
\usepackage{parskip} % Disable initial indent
\usepackage{attachfile}

\usepackage[utf8]{inputenc}
\fancyhf{}
\renewcommand{\headrulewidth}{0pt} % Remove default underline from header package
\rhead{CYBR/CMSC 691 - Lab 6}
%\rhead{}
\lhead{\begin{picture}(0,0) \put(0,-10){\includegraphics[width=1.1cm]{Images/UMBC-vertical}} \end{picture}}
\cfoot{\thepage}
\rfoot{Spring 2020}
\lfoot{Lab 6}
%\AtEndDocument{\vfill \hfill \LaTeX}

\begin{document}

\section*{Lab 6 - Data Pre-Processing}
\paragraph{}Using the provided Python Notebook, explore how data pre-processing can help or hinder model development. Experiment with various pre-processing functions provided by \texttt{sklearn} to see how model performance changes.

Documentation for \texttt{sklearn} transforms and pre-processing is available here: \url{https://scikit-learn.org/stable/data_transforms.html}.

\textbf{Points:} 20

\section*{Questions}
\paragraph{}In a separate document, provide your answers to the following questions. The expected format is a Word document, or similar, with a few images demonstrating your progress. The assignment will be submitted via Blackboard.

\begin{enumerate}
    \item Which pre-processing functions did you try, and how did they alter performance?
    \item Did the number of input features to the model change?
    \item What are some examples of how file features can be normalized to help improve model accuracy? Give specific examples.
\end{enumerate}

\end{document}